% GENERATED FILE. Edit canonical lessons, then run npm run generate:books.
% canonical-source-hash: fnv1a64:4e8320fd1aeffff9
% canonical-lessons: AR-C359-intiba, AR-C359-manzur, AR-C359-iftarada, AR-C359-iqtaraha, AR-C359-muqni

\chapter{Points of view}
\label{ch:ar-points-of-view}
\hlchaptermodality{359}

\begin{chapteropening}
\textbf{By the end of this chapter:} \emph{I can talk about an impression and a perspective, and say what someone assumed, suggested or found convincing.}

Complete the last lesson of chapter 359: I can talk about an impression and a perspective, and say what someone assumed, suggested or found convincing.
\end{chapteropening}

\section[\texorpdfstring{\ar{انطباع} (inṭibāʿ)}{انطباع (inṭibāʿ)}]{\ar{انطباع} (inṭibāʿ) — an impression}

\label{lesson:AR-C359-intiba}



\begin{quote}
Before the new one: say the Arabic for to prove, then the Arabic for to convince.
\end{quote}

\subsection*{You'll want to know: \ar{انطباع}}
\textbf{\ar{انطباع}} — \emph{inṭibāʿ} — \textquotedblleft{}an impression\textquotedblright{}.

\begin{quote}
\ar{ترك} \ar{الفندق} \ar{انطباعا} \ar{جيدا} \ar{عندي}.
\end{quote}

\begin{quote}
\emph{taraka al-funduq inṭibāʿan jayyidan ʿindī.}
\end{quote}

\begin{quote}
The hotel left a good impression on me.
\end{quote}

\begin{cousinweb}
\textbf{\ar{انطباع}} (\emph{inṭibāʿ}) is built on the native Arabic root \emph{ṭ-b-ʿ}, \textquotedblleft{}to stamp, to imprint\textquotedblright{} (Lane). \emph{Inṭibāʿ} is the Form VII verbal noun, being stamped: the mark a thing leaves on the mind, an impression.
\end{cousinweb}

\begin{grammarlens}[title={What you've built}]
One more word for giving reasons and opinions.
\end{grammarlens}

\subsection*{Your turn}
\begin{itemize}
\raggedright
  \item \emph{Say it:} \emph{inṭibāʿ}
  \item \emph{Say it:} \emph{inṭibāʿ}, once more
  \item \emph{Say it:} \emph{taraka al-funduq inṭibāʿan jayyidan ʿindī.} — The hotel left a good impression on me.
  \item \emph{Recall:} say the Arabic for to prove, then the Arabic for to convince, then say \emph{inṭibāʿ} again
\end{itemize}

\begin{tcolorbox}[breakable,colback=teal!4,colframe=teal!35!black,title={Before you move on}]
What does \ar{انطباع} mean? (\textquotedblleft{}An impression\textquotedblright{}.) Say it once more.
\end{tcolorbox}
\section[\texorpdfstring{\ar{منظور} (manẓūr)}{منظور (manẓūr)}]{\ar{منظور} (manẓūr) — a perspective}

\label{lesson:AR-C359-manzur}



\begin{quote}
Before the new one: say the Arabic for to convince, then the Arabic for an impression.
\end{quote}

\subsection*{You'll want to know: \ar{منظور}}
\textbf{\ar{منظور}} — \emph{manẓūr} — \textquotedblleft{}a perspective\textquotedblright{}.

\begin{quote}
\ar{من} \ar{منظور} \ar{الطبيب،} \ar{النوم} \ar{مهم} \ar{جدا}.
\end{quote}

\begin{quote}
\emph{min manẓūr aṭ-ṭabīb, an-nawm muhimm jiddan.}
\end{quote}

\begin{quote}
From the doctor's perspective, sleep is very important.
\end{quote}

\begin{cousinweb}
\textbf{\ar{منظور}} (\emph{manẓūr}) is built on the native Arabic root \emph{n-ẓ-r}, \textquotedblleft{}to look, to regard\textquotedblright{} (Lane). \emph{Manẓūr} is the passive participle, what is looked at; in modern use it is the angle you look from.
\end{cousinweb}

\begin{grammarlens}[title={What you've built}]
One more word for giving reasons and opinions.
\end{grammarlens}

\subsection*{Your turn}
\begin{itemize}
\raggedright
  \item \emph{Say it:} \emph{manẓūr}
  \item \emph{Say it:} \emph{manẓūr}, once more
  \item \emph{Say it:} \emph{min manẓūr aṭ-ṭabīb, an-nawm muhimm jiddan.} — From the doctor's perspective, sleep is very important.
  \item \emph{Recall:} say the Arabic for to convince, then the Arabic for an impression, then say \emph{manẓūr} again
\end{itemize}

\begin{tcolorbox}[breakable,colback=teal!4,colframe=teal!35!black,title={Before you move on}]
What does \ar{منظور} mean? (\textquotedblleft{}A perspective\textquotedblright{}.) Say it once more.
\end{tcolorbox}
\section[\texorpdfstring{\ar{افترض} (iftaraḍa)}{افترض (iftaraḍa)}]{\ar{افترض} (iftaraḍa) — to assume}

\label{lesson:AR-C359-iftarada}



\begin{quote}
Before the new one: say the Arabic for an impression, then the Arabic for a perspective.
\end{quote}

\subsection*{You'll want to know: \ar{افترض}}
\textbf{\ar{افترض}} — \emph{iftaraḍa} — \textquotedblleft{}to assume\textquotedblright{}.

\begin{quote}
\ar{افترضنا} \ar{أن} \ar{القطار} \ar{سيصل} \ar{مبكرا}.
\end{quote}

\begin{quote}
\emph{iftaraḍnā anna al-qiṭār sa-yaṣilu mubakkiran.}
\end{quote}

\begin{quote}
We assumed that the train would arrive early.
\end{quote}

Present: \ar{يفترض}, \emph{yaftariḍ}.

Followed by \ar{أن}: \ar{افترض} \ar{أن} …, assumed that ….

\begin{cousinweb}
\textbf{\ar{افترض}} (\emph{iftaraḍa}) is built on the native Arabic root \emph{f-r-ḍ}, \textquotedblleft{}to make a notch; to make binding\textquotedblright{} (Lane). It is Form VIII: to take a thing as laid down, to assume it.
\end{cousinweb}

\begin{grammarlens}[title={What you've built}]
One more word for giving reasons and opinions.
\end{grammarlens}

\subsection*{Your turn}
\begin{itemize}
\raggedright
  \item \emph{Say it:} \emph{iftaraḍa}
  \item \emph{Say it:} \emph{iftaraḍa}, once more
  \item \emph{Say it:} \emph{iftaraḍnā anna al-qiṭār sa-yaṣilu mubakkiran.} — We assumed that the train would arrive early.
  \item \emph{Recall:} say the Arabic for an impression, then the Arabic for a perspective, then say \emph{iftaraḍa} again
\end{itemize}

\begin{tcolorbox}[breakable,colback=teal!4,colframe=teal!35!black,title={Before you move on}]
What does \ar{افترض} mean? (\textquotedblleft{}To assume\textquotedblright{}.) Say it once more.
\end{tcolorbox}
\section[\texorpdfstring{\ar{اقترح} (iqtaraḥa)}{اقترح (iqtaraḥa)}]{\ar{اقترح} (iqtaraḥa) — to suggest}

\label{lesson:AR-C359-iqtaraha}



\begin{quote}
Before the new one: say the Arabic for a perspective, then the Arabic for to assume.
\end{quote}

\subsection*{You'll want to know: \ar{اقترح}}
\textbf{\ar{اقترح}} — \emph{iqtaraḥa} — \textquotedblleft{}to suggest\textquotedblright{}.

\begin{quote}
\ar{اقترح} \ar{المعلم} \ar{رحلة} \ar{إلى} \ar{المتحف}.
\end{quote}

\begin{quote}
\emph{iqtaraḥa al-muʿallim riḥla ilā al-matḥaf.}
\end{quote}

\begin{quote}
The teacher suggested a trip to the museum.
\end{quote}

Present: \ar{يقترح}, \emph{yaqtariḥ}.

The noun \ar{اقتراح}, \emph{iqtirāḥ}, a suggestion, comes from it.

\begin{cousinweb}
\textbf{\ar{اقترح}} (\emph{iqtaraḥa}) is built on the native Arabic root \emph{q-r-ḥ}, \textquotedblleft{}to wound\textquotedblright{} (Lane). It is Form VIII. Its modern sense, to put an idea forward, is usually traced to an older one, to bring out something new.
\end{cousinweb}

\begin{grammarlens}[title={What you've built}]
One more word for giving reasons and opinions.
\end{grammarlens}

\subsection*{Your turn}
\begin{itemize}
\raggedright
  \item \emph{Say it:} \emph{iqtaraḥa}
  \item \emph{Say it:} \emph{iqtaraḥa}, once more
  \item \emph{Say it:} \emph{iqtaraḥa al-muʿallim riḥla ilā al-matḥaf.} — The teacher suggested a trip to the museum.
  \item \emph{Recall:} say the Arabic for a perspective, then the Arabic for to assume, then say \emph{iqtaraḥa} again
\end{itemize}

\begin{tcolorbox}[breakable,colback=teal!4,colframe=teal!35!black,title={Before you move on}]
What does \ar{اقترح} mean? (\textquotedblleft{}To suggest\textquotedblright{}.) Say it once more.
\end{tcolorbox}
\section[\texorpdfstring{\ar{مقنع} (muqniʿ)}{مقنع (muqniʿ)}]{\ar{مقنع} (muqniʿ) — convincing}

\label{lesson:AR-C359-muqni}



\begin{quote}
Before the new one: say the Arabic for to assume, then the Arabic for to suggest.
\end{quote}

\subsection*{You'll want to know: \ar{مقنع}}
\textbf{\ar{مقنع}} — \emph{muqniʿ} — \textquotedblleft{}convincing\textquotedblright{}.

\begin{quote}
\ar{كان} \ar{كلام} \ar{المدير} \ar{مقنعا}.
\end{quote}

\begin{quote}
\emph{kāna kalām al-mudīr muqniʿan.}
\end{quote}

\begin{quote}
The manager's words were convincing.
\end{quote}

After \ar{كان} the adjective takes -an: \ar{كان} \ar{مقنعا}.

\begin{cousinweb}
\textbf{\ar{مقنع}} (\emph{muqniʿ}) is built on the native Arabic root \emph{q-n-ʿ}, \textquotedblleft{}to be content, satisfied\textquotedblright{} (Lane). \emph{Muqniʿ} is the Form IV active participle: that which convinces.
\end{cousinweb}

\begin{grammarlens}[title={What you've built}]
One more word for giving reasons and opinions.
\end{grammarlens}

\subsection*{Your turn}
\begin{itemize}
\raggedright
  \item \emph{Say it:} \emph{muqniʿ}
  \item \emph{Say it:} \emph{muqniʿ}, once more
  \item \emph{Say it:} \emph{kāna kalām al-mudīr muqniʿan.} — The manager's words were convincing.
  \item \emph{Recall:} say the Arabic for to assume, then the Arabic for to suggest, then say \emph{muqniʿ} again
\end{itemize}

\begin{tcolorbox}[breakable,colback=teal!4,colframe=teal!35!black,title={Before you move on}]
What does \ar{مقنع} mean? (\textquotedblleft{}Convincing\textquotedblright{}.) Say it once more.
\end{tcolorbox}
